% GENERATED FILE. Edit canonical lessons, then run npm run generate:books.
% canonical-source-hash: fnv1a64:616599ccfb19cc68
% canonical-lessons: MW-C128-pagri, MW-C128-ghaghro, MW-C128-lugro, MW-C128-kachali, MW-C128-angarkhi

\chapter{Turban, Long skirt, Woman's wrap, Bodice, Long coat}
\label{ch:mw-turban-long-skirt-woman-s-wrap-bodice-long-coat}
\hlchaptermodality{128}

\begin{chapteropening}
\textbf{By the end of this chapter:} \emph{I can say a turban, a long skirt, a woman's wrap, a bodice and a long coat.}

Complete the last lesson of chapter 128: I can say a turban, a long skirt, a woman's wrap, a bodice and a long coat.
\end{chapteropening}

\section[\texorpdfstring{\mw{पगड़ी} (pagṛī)}{पगड़ी (pagṛī)}]{\mw{पगड़ी} (pagṛī) — a turban}

\label{lesson:MW-C128-pagri}



\begin{quote}
Before the new one: say the Marwadi for a plough, then the Marwadi for a bullock cart.
\end{quote}

\subsection*{You'll want to know: \mw{पगड़ी}}
\textbf{\mw{पगड़ी}} — \emph{pagṛī} — \textquotedblleft{}a turban\textquotedblright{}.

\begin{cousinweb}
\textbf{\mw{पगड़ी}} (\emph{pagṛī}) is a tadbhava, a word inherited and reshaped by sound change. It comes from \emph{*paggā}, \textquotedblleft{}headdress\textquotedblright{}, a form CDIAL reconstructs from its descendants. CDIAL gives no Marwari form in this entry (CDIAL 7644); the match is through Hindi \emph{pagṛī}. CDIAL thinks a source in Sanskrit \emph{pragraha} scarcely likely, and explains the \emph{ṛ} as an added suffix.
\end{cousinweb}

\begin{grammarlens}[title={What you've built}]
One more word for this chapter.
\end{grammarlens}

\subsection*{Your turn}
\begin{itemize}
\raggedright
  \item \emph{Say it:} \emph{pagṛī}
  \item \emph{Say it:} \emph{pagṛī}, once more
  \item \emph{Say it:} say \emph{pagṛī}
  \item \emph{Recall:} say the Marwadi for a plough, then the Marwadi for a bullock cart, then say \emph{pagṛī} again
\end{itemize}

\begin{tcolorbox}[breakable,colback=teal!4,colframe=teal!35!black,title={Before you move on}]
What does \mw{पगड़ी} mean? (\textquotedblleft{}A turban\textquotedblright{}.) Say it once more.
\end{tcolorbox}
\section[\texorpdfstring{\mw{घाघरो} (ghāghro)}{घाघरो (ghāghro)}]{\mw{घाघरो} (ghāghro) — a long skirt}

\label{lesson:MW-C128-ghaghro}



\begin{quote}
Before the new one: say the Marwadi for a bullock cart, then the Marwadi for a turban.
\end{quote}

\subsection*{You'll want to know: \mw{घाघरो}}
\textbf{\mw{घाघरो}} — \emph{ghāghro} — \textquotedblleft{}a long skirt\textquotedblright{}.

\begin{cousinweb}
\textbf{\mw{घाघरो}} (\emph{ghāghro}) is a tadbhava, a word inherited from Sanskrit and reshaped by sound change. It comes from \emph{ghargharī}, \textquotedblleft{}girdle of small bells\textquotedblright{}. CDIAL gives no Marwari form in this entry (CDIAL 4444); the match is through Hindi \emph{ghāghrā} and Gujarati \emph{ghāghro}.
\end{cousinweb}

\begin{grammarlens}[title={What you've built}]
One more word for this chapter.
\end{grammarlens}

\subsection*{Your turn}
\begin{itemize}
\raggedright
  \item \emph{Say it:} \emph{ghāghro}
  \item \emph{Say it:} \emph{ghāghro}, once more
  \item \emph{Say it:} say \emph{ghāghro}
  \item \emph{Recall:} say the Marwadi for a bullock cart, then the Marwadi for a turban, then say \emph{ghāghro} again
\end{itemize}

\begin{tcolorbox}[breakable,colback=teal!4,colframe=teal!35!black,title={Before you move on}]
What does \mw{घाघरो} mean? (\textquotedblleft{}A long skirt\textquotedblright{}.) Say it once more.
\end{tcolorbox}
\section[\texorpdfstring{\mw{लूगड़ो} (lūgṛo)}{लूगड़ो (lūgṛo)}]{\mw{लूगड़ो} (lūgṛo) — a woman's wrap}

\label{lesson:MW-C128-lugro}



\begin{quote}
Before the new one: say the Marwadi for a turban, then the Marwadi for a long skirt.
\end{quote}

\subsection*{You'll want to know: \mw{लूगड़ो}}
\textbf{\mw{लूगड़ो}} — \emph{lūgṛo} — \textquotedblleft{}a woman's wrap\textquotedblright{}.

\begin{cousinweb}
\textbf{\mw{लूगड़ो}} (\emph{lūgṛo}) is a tadbhava, a word inherited and reshaped by sound change. It comes from \emph{*lugga}, \textquotedblleft{}defective\textquotedblright{}, a form CDIAL reconstructs from its descendants. In this large group of \textquotedblleft{}defective\textquotedblright{} words, Hindi \emph{lugṛā} means \textquotedblleft{}rag\textquotedblright{} and Gujarati \emph{lugṛũ} \textquotedblleft{}clothes\textquotedblright{}, and CDIAL's additions give a Rajasthani \emph{lūgṛī}, \textquotedblleft{}scarf, a woman's coloured wrap\textquotedblright{} (CDIAL 11072). How the senses connect is not settled.
\end{cousinweb}

\begin{grammarlens}[title={What you've built}]
One more word for this chapter.
\end{grammarlens}

\subsection*{Your turn}
\begin{itemize}
\raggedright
  \item \emph{Say it:} \emph{lūgṛo}
  \item \emph{Say it:} \emph{lūgṛo}, once more
  \item \emph{Say it:} say \emph{lūgṛo}
  \item \emph{Recall:} say the Marwadi for a turban, then the Marwadi for a long skirt, then say \emph{lūgṛo} again
\end{itemize}

\begin{tcolorbox}[breakable,colback=teal!4,colframe=teal!35!black,title={Before you move on}]
What does \mw{लूगड़ो} mean? (\textquotedblleft{}A woman's wrap\textquotedblright{}.) Say it once more.
\end{tcolorbox}
\section[\texorpdfstring{\mw{कांचळी} (kā̃chaḷī)}{कांचळी (kā̃chaḷī)}]{\mw{कांचळी} (kā̃chaḷī) — a bodice}

\label{lesson:MW-C128-kachali}



\begin{quote}
Before the new one: say the Marwadi for a long skirt, then the Marwadi for a woman's wrap.
\end{quote}

\subsection*{You'll want to know: \mw{कांचळी}}
\textbf{\mw{कांचळी}} — \emph{kā̃chaḷī} — \textquotedblleft{}a bodice\textquotedblright{}.

\begin{cousinweb}
\textbf{\mw{कांचळी}} (\emph{kā̃chaḷī}) is a tadbhava, a word inherited from Sanskrit and reshaped by sound change. It comes from \emph{kañculī}, \textquotedblleft{}bodice\textquotedblright{}. CDIAL gives no Marwari form in this entry (CDIAL 2627); the match is through Gujarati \emph{kā̃cḷī}. Its base word, \emph{*kañcu}, \textquotedblleft{}snake's skin; bodice\textquotedblright{}, does have Old Marwari \emph{kaṃcū}, \textquotedblleft{}bodice\textquotedblright{} (CDIAL 2626).
\end{cousinweb}

\begin{grammarlens}[title={What you've built}]
One more word for this chapter.
\end{grammarlens}

\subsection*{Your turn}
\begin{itemize}
\raggedright
  \item \emph{Say it:} \emph{kā̃chaḷī}
  \item \emph{Say it:} \emph{kā̃chaḷī}, once more
  \item \emph{Say it:} say \emph{kā̃chaḷī}
  \item \emph{Recall:} say the Marwadi for a long skirt, then the Marwadi for a woman's wrap, then say \emph{kā̃chaḷī} again
\end{itemize}

\begin{tcolorbox}[breakable,colback=teal!4,colframe=teal!35!black,title={Before you move on}]
What does \mw{कांचळी} mean? (\textquotedblleft{}A bodice\textquotedblright{}.) Say it once more.
\end{tcolorbox}
\section[\texorpdfstring{\mw{अंगरखी} (aṅgarkhī)}{अंगरखी (aṅgarkhī)}]{\mw{अंगरखी} (aṅgarkhī) — a long coat}

\label{lesson:MW-C128-angarkhi}



\begin{quote}
Before the new one: say the Marwadi for a woman's wrap, then the Marwadi for a bodice.
\end{quote}

\subsection*{You'll want to know: \mw{अंगरखी}}
\textbf{\mw{अंगरखी}} — \emph{aṅgarkhī} — \textquotedblleft{}a long coat\textquotedblright{}.

\begin{cousinweb}
\textbf{\mw{अंगरखी}} (\emph{aṅgarkhī}) is a tadbhava, a word inherited and reshaped by sound change. It comes from \emph{*aṅgarakṣa}, \textquotedblleft{}garment\textquotedblright{}, a form CDIAL reconstructs from its descendants. CDIAL gives no Marwari form in this entry (CDIAL 120); the match is through Hindi \emph{ãgarkhī} and Gujarati \emph{ãgarkhī}. It is built on \emph{aṅga}, \textquotedblleft{}body\textquotedblright{}, and \emph{rakṣa}, \textquotedblleft{}guard\textquotedblright{}.
\end{cousinweb}

\begin{grammarlens}[title={What you've built}]
One more word for this chapter.
\end{grammarlens}

\subsection*{Your turn}
\begin{itemize}
\raggedright
  \item \emph{Say it:} \emph{aṅgarkhī}
  \item \emph{Say it:} \emph{aṅgarkhī}, once more
  \item \emph{Say it:} say \emph{aṅgarkhī}
  \item \emph{Recall:} say the Marwadi for a woman's wrap, then the Marwadi for a bodice, then say \emph{aṅgarkhī} again
\end{itemize}

\begin{tcolorbox}[breakable,colback=teal!4,colframe=teal!35!black,title={Before you move on}]
What does \mw{अंगरखी} mean? (\textquotedblleft{}A long coat\textquotedblright{}.) Say it once more.
\end{tcolorbox}
